\section{The exceptional small pair actions}
\label{sec:exceptional-pairs}
Put $\beta=\log_2(38/25)$ and $\gamma_5=\frac13\log_2 5$.
Throughout this section, a pair action means an actual transitive group
$U\leq C_2\wr S_s$ in its specified action on $2s$ points. Write
$E=U\cap C_2^s$, $T=U/E$, and $M_s=\F_2^s$ with its actual $T$ action.
Every normal subgroup mentioned below is a normal subgroup of this
original $U$, unless another group is specified.

\begin{lemma}[Relative socles and fixed prefixes]\label{lem:pair-relative}
Let $F\lhd U$ be a binary group and $S=U/F$. Suppose that, for every
$N\lhd U$, the multiplicity $m_X$ of each simple binary type $X$ in
$\Soc(U/N)\cap FN/N$ satisfies $m_X\leq q e_X$, where
$e_X=\dim_{\End_U X}X$. Suppose every quotient of $S$ has a faithful
tuple of at most $k$ irreducible characters. Then every $Q=U/N$ has a
faithful tuple consisting of a $k$-character prefix with joint binary
kernel, followed by $q$ general characters. Consequently
\[
 |\Epi(J,Q)|\leq |\Aut Q|[Q:O_2(Q)]^q
                   5^{kb/3}(38/25)^{qb}\qquad(J\leq S_b).
\]
If $Q$ is binary, the whole tuple can instead be counted in the known
binary-source mode. If a specified top quotient has a faithful tuple of
$r$ binary-primary linear characters, one may use those $r$ positions
and the same $q$ relative general positions.
For $F=E$, the relative hypothesis follows whenever every homogeneous
section $X^m$ of $M_s$ satisfies $m\leq q e_X$.
\end{lemma}
\begin{proof}
Put $A=FN/N$. A minimal normal subgroup of $Q$ contained in $A$ is
central in $A$ and elementary abelian. Thus $V=\Soc(Q)\cap A$ is a
semisimple binary $Q/A$-module. The dual $V^*$ has $q$ module generators:
for $X^m$ the required number is $\lceil m/e_X\rceil$, and distinct
simple types share positions. Induce each corresponding linear
character of $V$ to $Q$ and choose an irreducible constituent. Its
restriction to $V$ consists of conjugates of that linear character;
hence the intersection of these $q$ kernels with $V$ is trivial.
Inflate a faithful tuple of $Q/A$. The joint kernel of the resulting
tuple is normal in $Q$, lies in $A$, and meets $\Soc(Q)$ trivially,
so it is trivial. This does not extend a character to a nonabelian
$A$ or split the original extension. Apply
Proposition~\ref{prop:anchor-count} to the fixed prefix. When $F=E$,
$V$ is a section of the original $M_s$. All choices depend on $Q$,
not on the later complete source $J$.
\end{proof}

\begin{lemma}[Two elementary module bounds]\label{lem:small-monomial}
A transitive monomial module of dimension $r>1$ over a field has no
homogeneous section $X^m$ with $m>\lfloor r/2\rfloor$.
For the regular binary $V_4$ module
$W=\F_2[x,y]/(x^2,y^2)$, every submodule has head dimension at most two;
the unique submodule attaining two is the augmentation ideal
$I=(x,y)$, of dimension three. These statements also hold after a
field extension.
\end{lemma}
\begin{proof}
A transitive permutation group has a derangement: the average number
of fixed points is one, and the identity fixes more than one point.
Choose a monomial element whose permutation is a derangement. Its
module over the principal ideal domain $k[t]$ is a sum of at most
$\lfloor r/2\rfloor$ cyclic modules, one per cycle. A submodule of a
$c$-generated torsion module over a principal ideal domain is
$c$-generated: lift it to a submodule of a free module of rank $c$.
The same holds for a quotient. On the other hand $X^m$, restricted
to this element, requires at least $m$ generators, by any nonzero
primary component and the invariant-factor theorem.
For the second assertion a proper ideal lies in $I$, whose radical
is $I^2=\langle xy\rangle$. A one-dimensional image in $I/I^2$
generates a cyclic ideal, whereas a two-dimensional image forces
the entire $I$. The whole ring has cyclic head. This proves the
claim and its equality assertion over any field of characteristic two.
\end{proof}

\subsection{The finite inputs and what they certify}
\label{subsec:pair-finite-input}
We use the following finite input on \emph{top actions}, separately
from the universal lifting arguments that follow. For a simple binary
$T$-module $X$, write $\delta_X=\dim_{\F_2}X$. Define
\[
 \tau(T)=\max_{B\leq M_s,\ X}
 \left\lceil\frac{\dim_{\F_2}\Hom_T(B,X)}{\delta_X}\right\rceil,
 \qquad \kappa(T)=\max_{L\lhd T}c(T/L),
\]
where $c$ denotes the least size of a faithful irreducible character
tuple. The formula for $\tau$ is exactly the largest Schur demand of a
homogeneous section: the multiplicity of $X$ in the head of $B$ is
$\dim\Hom_T(B,X)/\dim\End_T(X)$.

\begin{lemma}[Certified top data]\label{lem:small-pair-data}
The top-action certificates have the following properties.
\begin{enumerate}[label=(\roman*)]
\item Among the $301$ transitive degree-twelve actions, $283$ have
$(\tau,\kappa)$ equal to
\[
 (1,1),(1,2),(1,3),(2,1),(2,2),(3,1),
\]
with respective multiplicities $71,102,19,33,56,2$.
The remaining eighteen consist of eight binary-core actions with
quotient $C_3$ or $S_3$ as in Lemma~\ref{lem:twelve-core}; eight actions
with an elementary normal ternary group on four triples as in
Lemma~\ref{lem:four-triple-compression}; and two paired actions with
six-point tops as in Lemma~\ref{lem:paired-six}.
\item For the two six-point groups
\[
 S_{36}=\langle(124),(24)(56),(13)(25)(46)\rangle,
 \quad
 S_{72}=\langle(124),(24),(13)(25)(46)\rangle,
\]
all homogeneous sections of $M_6$ have Schur demand at most one.
Their only simple types have dimensions one and four, with commuting
field $\F_2$. Every nontrivial quotient has a faithful irreducible
character except $C_2^2$, which has two faithful binary linear positions.
\item For all $24$ nonbinary transitive degree-eight groups $S$, every
nontrivial homogeneous type in $M_8$ has Schur demand at most one,
every trivial section has dimension at most two, and $\kappa(S)\leq2$.
If an invariant submodule has trivial head dimension two, it is a
unique five-dimensional submodule $B$. Every invariant $D\geq B$
satisfies $\dim D/(B+[D,S])\leq1$.
\item The $29$ zero-ternary minimal-four degree-sixteen tops have
$\tau\leq4$, $\kappa\leq2$. The $15$ degree-sixteen tops with two
primitive affine eight-point blocks have $\kappa\leq1$ and binary
permutation composition Schur demand at most four. The $20$ primitive
affine degree-sixteen tops have the respective bounds two and four.
The two inequivalent nontrivial binary composition factors of the
projective eight-point $\operatorname{PSL}_2(7)$ module have dimension
three, commuting field $\F_2$, and each occurs once; the trivial factor
occurs twice.
\end{enumerate}
\end{lemma}
\begin{proof}
These are the finite calculations recorded by the pair-top certificate
contract in Appendix~\ref{sec:certificates}. We specify their mathematical
coverage, since a list of tested module examples would not suffice.
Starting with every vector, form its cyclic submodule and close the
resulting list under sums. Every invariant subspace is the sum of
cyclic submodules of its vectors, so this constructs the complete
invariant-subspace lattice. Maximal intervals give the simple modules;
intertwining linear equations on the same generators identify their
isomorphism types, their endomorphism fields, and the displayed Hom
spaces. A composition chain gives the separate composition bounds.
Thus every numerator and every homogeneous quotient is covered.
For normal subgroups start with kernels of all irreducible characters
and close under intersections. Every normal subgroup occurs, since
the regular representation of its quotient is faithful and decomposes
into irreducibles. Literal kernels, not quotient isomorphism types,
are retained. Searching these kernel intersections proves the asserted
covers for all quotients.

Completeness of the degree-twelve and primitive action lists is the
published classification input
\cite{Hulpke2005,HoltRoyle2020,RoneyDougal2005Primitive,TransGrp365,PrimGrp344}.
The imprimitive degree-eight and two-affine-block constructions use
all normal quotients of the local group, all graph isomorphisms between
these quotients, and all block-swap cosets in their actual normalizers;
Goursat's lemma and the order-two coset condition prove completeness.
The four-pair constructions use the full invariant-kernel list and all
lifts of fixed top generators. The zero-ternary list similarly uses
all degree-four local subgroups, both induced binary linear frames,
all invariant kernels, and every solution of the generator-relator
cocycle equations. The computed cocycle-space dimension, modulo the
coboundary-space dimension, verifies that every complement class occurs.
Transporters identify actions in the stated lists. These procedures
certify only the finite top data. In particular they make no assertion
about unlisted normal subgroups of a subsequent pair lift.
\end{proof}

\subsection{The twelve-pair actions}

\begin{lemma}[The binary-core twelve-point tops]\label{lem:twelve-core}
Let $T$ be one of the eight binary-core actions in
Lemma~\ref{lem:small-pair-data}(i). Put $P=O_2(T)$.
For every pair lift $U$ and every $N\lhd U$, put $Q=U/N$.
If $T/P=C_3$, $Q$ admits one binary-primary linear and four general
faithful positions. If $Q$ is nonbinary, the linear position can be
made an ordinary linear anchor with binary kernel.
If $T/P=S_3$, every nonbinary $Q$ admits four general faithful positions,
one of which has binary kernel. Every binary $Q$ admits one
binary-primary linear and four general faithful positions.
\end{lemma}
\begin{proof}
First suppose $T/P=C_3$ and choose $h$ of order three. There are three
$P$-orbits of size four, and $h$ cycles them. Identify them by $h$.
Then $P\leq L^3$ has full local projections, with
$L=C_4,V_4$, or $D_8$. The simple binary types are $1$ and the
natural $C_3$ type $X$ of binary dimension two and Schur dimension one.
If $L=C_4$ or $D_8$, restrict each coordinate module to an actual
regular $C_4$. Its module is $\F_2[t]/(t^4)$, with cyclic submodules.
A three-coordinate filtration bounds every $P$-head by three; hence
the homogeneous-section bounds are three for $1$ and one for $X$.
If $L=V_4$, Lemma~\ref{lem:small-monomial} bounds a $P$-head by six.
Equality forces dimension three in each coordinate filtration and
rules out a full coordinate projection, since that would contribute
only one generator. Thus the numerator is the sum of the three
augmentation ideals. Its head, with $h$ cycling coordinates, is two
regular $C_3$ modules, namely $1^2\oplus X^2$. It cannot contain
$X^3$. Thus the $X$ bound is two.

In this last case a trivial $T$-section of rank four is also impossible.
Odd-order averaging $e=1+h+h^2$ would give, for its numerator $B$,
$B^h=M_{12}^h$ and $([B,P])^h=0$. Hence $e(p-1)e=0$ for every
$p=(p_1,p_2,p_3)\in P$, or on the regular four-point coordinate
\[
 R_{p_1}+R_{p_2}+R_{p_3}=I.
\]
The four regular permutation matrices are linearly independent.
Therefore $p=0$ or $p$ is a rotation of $(x,x,0)$ for a nonzero
$x\in V_4$. Rotation invariance and addition show that two independent
values of $x$ would produce a tuple with three nonzero entries, a
contradiction. The coordinate projections could not then be all $V_4$.
The trivial section bound is three.

Coprime centralizer lifting maps $C_P(h)$ onto the fixed subgroup in
every quotient of $P$ by a $T$-normal subgroup. Evaluation on one
coordinate embeds $C_P(h)$ in $L$. Consequently every quotient of
$T$ has central binary rank at most two. Moreover
$T'=[P,h]P'$ and the fixed part of $T'/P'$ is trivial, so
$C_{T'}(h)=C_{P'}(h)$; evaluation embeds this in $L'$, of order at
most two. Every quotient of $T$ therefore has derived central binary
rank at most one. Its nontrivial binary-type multiplicity is at most
two for $L=V_4$, and at most three otherwise: filter any elementary
section of $P\leq L^3$ through the coordinates, each of rank at most two.

For the whole $Q$, filter its binary socle by its intersection with
$EN/N$ and its image in $Q/(EN/N)$. Both are modules for the same
original $T$; same-type multiplicities add. The resulting central
rank is at most $3+2=5$, the derived central rank at most $3+1=4$,
and the nontrivial multiplicity at most $2+2=4$ or $1+3=4$.
The odd socle is central of rank at most one and has zero intersection
with $Q'$. There is no nonabelian socle. Theorem~\ref{thm:character-feasibility}
gives the simultaneous profile $(r,l,t)=(1,0,4)$.
If $Q$ is nonbinary, multiply the designated binary-primary linear
character by a faithful cubic character of $Q/O_2(Q)$. Its kernel
shrinks to its intersection with $O_2(Q)$, since a binary-power root
and a nontrivial cube root cannot have product one. This gives the
claimed ordinary linear anchor.

Now suppose $T/P=S_3$. The certified actions have either six $P$-orbits
of size two, with regular $S_3$ action on them, or three of size four,
with natural $S_3$ action. For $h$ of order three, $T_0=P\langle h\rangle$
has respectively two six-point orbits or one twelve-point orbit.
The simple binary types are $1$ and the natural $S_3$ type $Y$, whose
Schur dimension is two; restriction to $\langle h\rangle$ preserves
its multiplicity. A nonbinary $Q$ has binary kernel $N$, since the
odd part of $|U|$ is three. Its top quotient is $T/L$ for $L\leq P$.
Any elementary section of the actual $P\leq S_{12}$ has rank at most
six by \cite[Theorem 2.8]{RoneyDougalTracey2025}; its $Y$ multiplicity
is therefore at most three.

In the size-two case each local module is a cyclic regular $C_2$
module. Thus the pair-module $Y$ multiplicity is at most three.
Frattini's argument gives $T=P N_T(\langle h\rangle)$; take a binary
element in this normalizer mapping to a reflection. It swaps the two
triples of $P$-orbits and therefore acts without fixed points on the
four $h$-orbits. It has at most two cycles on these four labels.
Exact $h$-invariants and the cyclic module argument bound every trivial
$T$-section by two. In the size-four case the preceding $C_3$ proof
bounds pair-module trivial multiplicity by three and natural
multiplicity by two.

In both cases $C=C_P(h)$ acts faithfully on the four $h$-orbits:
the kernel in the ambient centralizer is odd. Thus $C\leq S_4$ is
binary, and every current top quotient has central binary rank at most
two by coprime fixed lifting. If that rank is two, $C$ is noncyclic
and contains a $V_4$. A faithful $V_4$ action on four points is regular
or has two two-point orbits with distinct kernels; in either case its
trivial module-section bound is two. Exact $h$-invariants therefore
improve the pair bound to two in this same case. The whole central
rank is at most four, and the whole natural Schur demand is at most
$\lceil(3+3)/2\rceil=3$. The odd socle rank is at most one.
Theorem~\ref{thm:character-feasibility} gives four general positions.
One has binary kernel: otherwise all four normal kernels contain a
Sylow three subgroup, hence all its conjugates, contradicting their
trivial intersection.

It remains to handle binary $Q$ in the $S_3$ case, for which the last
argument about the top quotient was not asserted. An actual Sylow
three element of $U$ is semiregular on the original twenty-four points,
with eight three-cycles, and is killed by $N$. Frattini gives
$U=N N_U(\langle h\rangle)$. In the ambient normalizer
$C_3^8:(C_2\times S_8)$, the common inversion factor $C_2$ must be
retained. The odd base of the actual normalizer is killed by $N$;
its image $H\leq C_2\times S_8$ is binary and surjects onto $Q$.
It has a faithful ten-point action, so all elementary sections have
rank at most five by the same published abelianization bound.
For a Sylow subgroup $P_8=C_2^4:D_8$ of $S_8$, one has
$|P_8'|=2^4$: the base commutator has order eight and $|D_8'|=2$.
Hence $|Q'|\leq16$. The central rank of $Q$ is at most five and its
derived central rank at most four. The simultaneous feasibility
criterion again gives one binary-primary linear and four general
positions, now in the known binary-source mode.
\end{proof}

\begin{lemma}[Four ternary orbits]\label{lem:four-triple-compression}
Suppose $T$ has an elementary normal ternary subgroup $C$ of rank at
least two with four orbits of size three, $T/C$ binary, and
$|O_2(T)|\leq2$. Every pair lift $U$ has a quotient $R$ with a
specified faithful twenty-point action and a constant $D_U$ such that
\[
 Z_J(U)\leq Z_J(R)+D_U5^b\qquad(J\leq S_b).
\]
\end{lemma}
\begin{proof}
Averaging over $C$ gives the $T$-equivariant decomposition
$M_{12}=M_0\oplus W$, with $\dim M_0=4$, $\dim W=8$.
The map $(v,t)\mapsto(\operatorname{avg}_C(v),t)$ on the ambient
semidirect product has kernel $W$. Restricted to $U$, its kernel is
$E_1=E\cap W=[E,C]$; put $R=U/E_1$. The ambient $M_0:T$ acts
faithfully on the disjoint union of the twelve pair labels and the
eight $C$-orbit labels on the original twenty-four points. The first
action detects $T$ and the second detects $M_0$, proving faithfulness
of this twenty-point representation after restriction to $R$.

Fix an original $N$ not containing $E_1$, and choose a minimal normal
$S$ in the nonzero image of $E_1$ in $Q=U/N$. On restriction to $C$,
$W$ consists of four nontrivial two-dimensional local types. Their
kernels have trivial intersection, and they form one $T$-orbit.
There are at least two distinct types, since $C$ has rank at least two;
each has multiplicity at most two. Every nonzero $T$-section of $W$
therefore acts faithfully on $C$. In particular if $L$ is the kernel
of the action of $T$ on $S$, then $L\cap C=1$, so $L$ is normal binary
and $L\leq O_2(T)$.

Put $A=EN/N$. Since $A$ is abelian, there is an exact sequence
\[
 1\longrightarrow A\longrightarrow C_Q(S)
 \longrightarrow L/\pi(N)\longrightarrow1.
\]
Every minimal normal subgroup of $Q$ commutes with $S$, including
$S$ itself. Thus the entire socle lies in this binary centralizer.
Its image modulo $A$ has dimension at most one and is trivial.
The trivial part in $A$ is a section of $M_0$, the ordinary permutation
module on four blocks, so has rank at most two by
Lemma~\ref{lem:small-monomial}. Every nontrivial type in $A$ lies in
$W$ and has multiplicity at most two by restriction to $C$.
The whole socle therefore has Schur demand at most three, and
Theorem~\ref{thm:character-feasibility} supplies three general positions.
Its epimorphisms cost at most $|\Aut Q|5^b$.
Finally the interval $N\geq E_1$ is exactly the normal-subgroup lattice
of $R$, with the same quotients and epimorphism sets. Summing that
interval gives exactly $Z_J(R)$; summing the other normals gives
the asserted constant $D_U=\sum_{N\not\supseteq E_1}|\Aut(U/N)|$.
\end{proof}

\begin{lemma}[The paired six-point tops]\label{lem:paired-six}
Suppose $U\to T\to S$ has degrees $24,12,6$, with both successive
systems consisting of pairs, and $S=S_{36}$ or $S_{72}$ from
Lemma~\ref{lem:small-pair-data}. Every original quotient admits either
one general anchor and three relative general positions, or two
binary-primary linears and three relative general positions in the
known binary-source mode. A binary quotient with no exceptional
$C_2^2$ top needs at most four general positions.
\end{lemma}
\begin{proof}
Let $F$ be the kernel of $U\to S$. On the six original four-point
blocks, $F\leq D_8^6$, and $F/E$ is a submodule of $M_6$.
Let $C_0\leq M_{12}$ be the vectors constant on the six pairs of
pair labels. Both $C_0$ and $M_{12}/C_0$ are $M_6$.
Fix $N$, put $A=FN/N$, and take a homogeneous component $X^m$ of
$\Soc(U/N)\cap A$. Its preimage $D\leq F$ has the three successive
factors obtained from
\[
 D\cap E\cap C_0,\qquad
 (D\cap E+C_0)/C_0,\qquad \pi(D)\leq F/E.
\]
The denominators are the respective images of $D\cap N$; the last
is $\pi(F\cap N)$, since $D$ contains $F\cap N$.
Thus all three factors are homogeneous sections of the same $M_6$
for the same original $S$ type. Their multiplicities add, and
Lemma~\ref{lem:small-pair-data}(ii) gives $m\leq3e_X$.
Apply Lemma~\ref{lem:pair-relative}. Except for top $C_2^2$, a
nontrivial top quotient has a one-character prefix with binary joint
kernel. For top $C_2^2$, $Q$ itself is binary and the prefix consists
of two full-order binary linear characters. A trivial top needs only
the three relative positions. No part of this proof assumes that
$F$ or $A$ is abelian.
\end{proof}

\begin{theorem}[Complete twelve-pair estimate]\label{thm:twelve-pairs}
The family of all actual pair actions with twelve pairs and nonbinary
top has a complete exponentially small forward row and a
quadratic-exponentially small scalar, with its original action weights.
\end{theorem}
\begin{proof}
For the $283$ generic tops, Lemma~\ref{lem:pair-relative} gives slope
$\kappa\gamma_5+\tau\beta$ and target-dependent finite constants.
For each of the six pairs in Lemma~\ref{lem:small-pair-data}(i),
\[
 5^{64\kappa}38^{192\tau}<2^{573}25^{192\tau},
\]
so the slope is less than $3-1/64$. The binary-core slopes are
$\frac13\log_2 3+4\beta$, $\frac12+4\beta$, and
$\gamma_5+3\beta$, all less than $3-1/64$. The paired-six slopes
are $\gamma_5+3\beta$, $1+3\beta$, and $4\beta$, with the same
strict margin. These inequalities follow, without rounding logarithms,
by raising to the $192$nd power; the corresponding numerators are
$3^{64}38^{768}$, $2^{96}38^{768}$, $5^{64}38^{576}$,
$2^{192}38^{576}$, and $38^{768}$, against $2^{573}$ times the
matching power of $25$. Proposition~\ref{prop:anchor-count} and
Corollary~\ref{cor:character-forward} sum every literal normal in these
finite original width-twenty-four actions. Constants such as
$[Q:O_2(Q)]^q|\Aut Q|$ are fixed before choosing $J$.
For the other eight tops, Lemma~\ref{lem:four-triple-compression}
and Theorem~\ref{thm:finite-comparator} apply with $v=20<24$ and
$\eta=\log_2 5<3$. This last assertion includes the entire normal
interval of the quotient, rather than only its faithful target.
There are finitely many subgroups of $S_{24}$, so the complete union
is covered by one fusion, retaining the factor
$1/(b!a(U)p_nG_{\lfloor n/2\rfloor})$ in each original hot sum.
\end{proof}
